\documentclass[10pt,a4paper]{amsart}
\usepackage[margin=19mm]{geometry}
\usepackage{amsmath,amssymb,mathtools}
\usepackage[scr=rsfs,cal=euler]{mathalfa}
\usepackage{xfrac}
\usepackage[unicode,hidelinks,bookmarks=false]{hyperref}
\newcommand{\ie}{i.e.\ }
\newcommand{\cf}{cf.\ }
\newcommand{\resp}{respectively\ }
\newcommand{\Subsection}{\subsection}
\providecommand{\nobreakdash}{\nobreak\hbox{-}}
\setlength{\emergencystretch}{2em}
\multlinegap0pt
\usepackage{amssymb}
\usepackage{mathtools}
\usepackage[shortlabels]{enumitem}
\usepackage{float}
\usepackage{xcolor}
\usepackage{tikz}
\newtheorem{proposition}{Proposition}
\title{\LARGE\bfseries Weak-Coupling Limit of the Lattice Nonlinear Schr\"odinger Integral Equation}
\author{Felipe Taha Sant'Ana}
\date{Source: arXiv:2603.09522v3 (CC BY 4.0)}
\begin{document}
\maketitle

\noindent\emph{Source and licence.} \LARGE\bfseries Weak-Coupling Limit of the Lattice Nonlinear Schr\"odinger Integral Equation. Version arXiv:2603.09522v3 (CC BY 4.0), distributed by arXiv under the Creative Commons Attribution 4.0 International licence, http://creativecommons.org/licenses/by/4.0/, as recorded in the arXiv licence metadata for this exact version and shown on the abstract page https://arxiv.org/abs/2603.09522v3. Full licence notice: \texttt{licenses/arxiv-2603.09522-CC-BY-4.0.txt}.\par\smallskip

\noindent\textit{Editorial scope.} Counted unit: the article's proposition prop:gap (source0.tex lines 2223--2236). The article's complete original proof of it is delivered with it (source0.tex lines 2238--2265, one \texttt{proof} environment of 28 lines). No mathematics is written, repaired or re-derived: the statement and the proof are the article's own lines, in the article's own notation, and no retained line is substituted or re-flowed.  Retained uncounted as the source-local context the proof needs: lines 559--563; lines 2131--2137.  Nothing was omitted from any retained span, so the package carries no omission record.  No retained line contains a citation, so the wrapper carries no bibliography and no bibliography entry.  No retained line contains a figure environment, an image-inclusion command or drawing code, so no image is delivered and the package declares no asset dependency.  Editorial formatting only, in addition: an AMS \texttt{amsart} preamble that restates the article's own declarations and macros used by the retained lines, namely the environments \texttt{proposition} on separate counters, together with the standard packages those lines need.  The retained environments print in this document as Proposition 1 in order of appearance, whereas the article prints the same items with the numbering of its own full document.  The complete native source of the article as distributed in the arXiv e-print for this exact version is bundled unchanged as \texttt{sources/196069/source0.tex}.
\begin{equation}\label{eq.K-fourier}
  \hat K(p;\,\kappa) = \int_{-\infty}^{\infty}e^{-ip\lambda}\,
  \frac{2\kappa}{\kappa^2 + \lambda^2}\,d\lambda
  = 2\pi\,e^{-\kappa|p|}\,.
\end{equation}

\begin{equation}\label{eq.minmax}
    \lambda_n(Q)
    \;=\;
    \max_{\substack{V \subset L^2([-Q,Q]) \\[1pt] \dim V = n+1}}
    \;\min_{f \in V,\; \|f\|=1}
    \;\langle f, \mathcal{K}_Q\, f \rangle
\end{equation}

\begin{proposition}[Two-sided bounds on the spectral gap]\label{prop:gap}
  For all $Q \geq \pi/4$,
  \begin{equation}\label{eq.gap-upper}
    \frac{\pi^2\left(1 - e^{-1}\right)}{4\,Q}
    \;\leq\;
    \Delta_0(Q)
    \;\leq\;
    \frac{2\pi\lambda_1}{Q}\,,
    \qquad \lambda_1 = 1.157774\ldots\,,
  \end{equation}
  where $\lambda_1$ is the smallest eigenvalue of the half-Laplacian
  $\mathcal{A} = (-d^2/d\xi^2)^{1/2}$ on $(-1,1)$ with exterior Dirichlet
  condition.
\end{proposition}

\begin{proof}
Both bounds are comparisons of symbols under the min-max
principle~\eqref{eq.minmax}, applied to
$\langle u, (2\pi I - \mathcal{K}_Q)u\rangle
= \int \sigma(p)\,|\hat u(p)|^2\,dp$ for $u$ supported in $[-Q,Q]$ with
$\|u\| = 1$, where $\sigma(p) = 1-e^{-|p|}$ and the Fourier normalisation is
that of~\eqref{eq.K-fourier}.

\emph{Upper bound.}  $1-e^{-x} \leq x$ for $x \geq 0$, so
$\sigma(p) \leq |p|$ pointwise and $2\pi I - \mathcal{K}_Q \leq 2\pi\mathcal{A}$
as quadratic forms on functions supported in $[-Q,Q]$.  The eigenvalues of
$\mathcal{A}$ on $(-Q,Q)$ are $\lambda_n/Q$, by the homogeneity
$\lambda_n(kD) = k^{-1}\lambda_n(D)$ of a symbol of degree one, so
$\Delta_0(Q) \leq 2\pi\lambda_1/Q$.

\emph{Lower bound.}  Since $(1-e^{-x})/x$ decreases,
$\sigma(p) \geq (1-e^{-1})\min(|p|,1)$ for all $p$.  For $u$ supported in
$[-Q,Q]$ with $\|u\|=1$, Cauchy--Schwarz gives $|\hat u(p)|^2 \leq 2Q$, hence
for any $\epsilon \leq 1$
\begin{equation}
  \frac{1}{2\pi}\int \min(|p|,1)\,|\hat u|^2\,dp
  \;\geq\; \epsilon\left[1 - \frac{1}{2\pi}\int_{|p|<\epsilon}|\hat u|^2 dp\right]
  \;\geq\; \epsilon\left[1 - \frac{2\epsilon Q}{\pi}\right]\,.
\end{equation}
Choosing $\epsilon = \pi/(4Q)$, which is admissible for $Q \geq \pi/4$, the
right-hand side is $\pi/(8Q)$, and multiplying by $2\pi(1-e^{-1})$ gives the
stated bound.
\end{proof}

\end{document}
